\section{Proofs for objective-cutoff propagation}
\label{app:propagation}

\subsection{Greatest fixed boxes and threshold attainment}
\label{app:propagation-fixed-point}

\begin{proof}[Proof of \cref{prop:fixed-point}]
An elementary revise is monotone and deflationary. Take the box hull $H$
of the union of all consistent boxes inside $Z_0(C)$, including the empty
box. Each such box $A$ satisfies
$A=\rho_E(A)\subseteq\rho_E(H)$. Thus $H\subseteq\rho_E(H)$, and
deflation gives equality. Its root interval is below $c$. This proves
existence and maximality. If a step starts from a box containing $H$,
monotonicity gives
$H=\rho_E(H)\subseteq\rho_E(Z)\subseteq\rho(Z)$, so every step preserves it.

Exact revises are continuous along decreasing compact boxes. To see this,
let $Y_k\downarrow Y$. The closed sets $E\cap Y_k$ are nested and compact.
If their intersection is empty, one is already empty. Otherwise maxima
and minima of every coordinate converge to the extrema on $E\cap Y$:
take extremizers, extract a convergent subsequence, and use closedness.
Consequently $\rho_E(Y)=\bigcap_k\rho_E(Y_k)$.
For a fair run $Z_k$, let $Z_\infty=\bigcap_kZ_k$. Along the infinitely
many applications of $E$, this identity implies
$Z_\infty\subseteq\rho_E(Z_\infty)$. Every constraint and the cutoff
therefore fix the limit. Preservation and maximality give $Z_\infty=H$.
An empty intersection of nested nonempty compact boxes is impossible,
so an empty limit is reached at a finite iterate.

A forward partial revise updates the parent projection of $E\cap Z$;
a backward partial revise updates its child projections. Each is monotone,
deflationary, continuous from above, and preserves consistent boxes.
A box fixed by both has all the projections required by the exact revise.
The same fair-limit argument therefore applies to these partial steps.

Shrinking $C$ or decreasing $c$ shrinks $Z^*(C,c)$, by maximality.
For a decreasing sequence $c_k\downarrow c$ with nonempty fixed boxes,
their intersection is nonempty and fixed at cutoff $c$. If all their
variable projections contain $y$, the intersection also contains $y$ in
its projection. Thus the two nonempty-cutoff sets are closed upper sets.
The exact lifted evaluation at $y$ is a singleton consistent box at
cutoff $f(y)$. No box is consistent below $F_{\rm lo}(C)$. This proves
attainment and \eqref{prop:bound-chain}. The closed upper sets are precisely
$[\pi_{\mathcal D}(C),\infty)$ and
$[\pi_{\mathcal D}(C,y),\infty)$; preservation and convergence prove the
strict emptiness and removal statements.
\end{proof}

We will also use a consequence of topological forward evaluation. Every
consistent box $Z$ with constants fixed satisfies
\begin{equation}
 Z\subseteq Z_0(\Pi_xZ).
 \label{prop:own-forward}
\end{equation}
Indeed variables and constants satisfy this inclusion first. Supporting
each operation-node endpoint shows that its interval is contained in its
forward operation image. Induction through the graph proves the claim.

\subsection{First-order loss and owned event accounting}
\label{app:propagation-transfer}

For a centered subbox $U=y+\prod_i[-\delta_i,\delta_i]$ define
$a_j=\sum_i|\partial_it_j(y)|\delta_i$. If every second partial of $t_j$
is bounded in absolute value by $H_j$, Taylor's theorem gives
\[
 \min_Ut_j\le t_j(y)-a_j+\frac{H_j}2\|\delta\|_1^2,
 \qquad
 \max_Ut_j-\min_Ut_j\le2a_j+H_j\|\delta\|_1^2.
\]
The first inequality uses the corner opposing the derivative signs;
the second bounds the two extreme values. The witness lemma yields
\begin{equation}
 \pi_{\mathcal D}(C,y)\le f(y)-
 \left(\sum_ja_j-2\max_ja_j\right)
 +\left(\frac12\sum_jH_j+\max_jH_j\right)\|\delta\|_1^2.
 \label{prop:taylor-loss}
\end{equation}
For a coordinate segment only the corresponding second partial is needed.

\begin{proof}[Point localization in \cref{prop:cnd-transfer}]
Put $g=m(y)+\epsilon$ and choose
$\delta=\min\{y_i-l_i,u_i-y_i,\rho_*\}$. By strict removal and
\eqref{prop:taylor-loss},
\[
 f_* -\epsilon<\pi_{\mathcal D}(C,y)
 \le f(y)-d_0\delta+K_0\delta^2
 \le f(y)-d_0\delta/2.
\]
Therefore $\delta<2g/d_0<\rho_*$, which gives the nearest-face bound.
Each coordinate product is at most $w(C)$ times its nearest-face
distance. Summing gives $q_C(y)<2nw(C)g/d_0$, hence
\eqref{prop:cnd-validity}. A degenerate box has $q_C=0$ in its
degenerate directions and causes no difficulty.
\end{proof}

For completeness we establish the propagation event statement with
inherited bounds. A phase starts with variable box $B_0$, ends with a
retained box $B_f$, and consists of consecutive objective-only runs.
Each restart uses the forward box of its current variable box, intersected
with final lifted boxes inherited from earlier runs at this node or its
ancestors. Its \emph{inheritance chain} includes all those earlier phases
recursively. Let $c_\Phi$ be the smallest cutoff in this chain. All cutoffs
are at least $f_* -\epsilon$, even if incumbent values change.

Every run in phase $\Phi$ preserves
$H_\Phi=Z^*(B_0,c_\Phi)$. Prove this in time order. An earlier phase
$\Psi$ feeding this phase has $B_0\subseteq B_{0,\Psi}$ and
$c_\Phi\le c_\Psi$, so $H_\Phi\subseteq H_\Psi$ by monotonicity.
Its inherited final box contains $H_\Phi$. Earlier runs in the same
phase also contain $H_\Phi$. Their variable projections contain
$\Pi_xH_\Phi$, so restarted forward boxes contain $H_\Phi$ by
\eqref{prop:own-forward}. Finally $H_\Phi$ is consistent for each larger
cutoff used in this phase. Preservation proves the induction. Thus
$\Pi_xH_\Phi\subseteq B_f$.

If $B_f=\prod_i[a_i,b_i]$ is nonempty, give each removed point to its first
coordinate outside $[a_i,b_i]$. The resulting owners have preceding
coordinates in their retained intervals, the first offending coordinate
in $[l_i,a_i)$ or $(b_i,u_i]$, and subsequent coordinates in their original
intervals. There are at most $2n$ convex rectangular owners. Their closures
are the test boxes. Carry the retained boundaries forward. For an owner
with test box $S\subseteq B_0$, monotonicity gives
$Z^*(S,c_\Phi)\subseteq H_\Phi$. Every owned $y\notin B_f$ therefore
satisfies $\pi_{\mathcal D}(S,y)>c_\Phi\ge f_* -\epsilon$.
If a phase empties its box, preservation implies $H_\Phi=\varnothing$;
its whole terminal owner satisfies this strict test. Splits similarly
assign their shared boundary to one child. Following each point through
the finite run partitions the domain into terminal and discarded owners,
each satisfying either the relaxation point test or the propagation test.
Extra variable-space reductions may be included if their discarded owners
have the stated relaxation point certificate; each such reduction adds
at most $2n$ pieces to the accounting.

Every point of $N\cap\{g<\vartheta\}$ consequently satisfies
$g\ge\alpha_{\rm eff}q_C$ in its own test box. If $g\le\epsilon+\eta$,
the nearest endpoint in each coordinate lies within
$r=\sqrt{(\epsilon+\eta)/\alpha_{\rm eff}}$, since a coordinate product
is at least the square of that nearest distance. The $2^n$ vertex balls
of each test box cover its certified points. This proves
\eqref{prop:cnd-cover}.

For a nondegenerate test box, arithmetic--geometric mean gives
\[
 \int_Cq_C(y)^{-n/2}\,dy
 \le n^{-n/2}\prod_{i=1}^n
 \int_{l_i}^{u_i}((y_i-l_i)(u_i-y_i))^{-1/2}\,dy_i
 =n^{-n/2}\pi^n.
\]
Thus its certified owner contributes at most
$(\pi^2/(\alpha_{\rm eff}n))^{n/2}$ to the integral. Degenerate boxes
have zero $n$-dimensional measure. Sum over owners to prove
\eqref{prop:cnd-integral}. A compact smooth $p$-dimensional minimizing
patch has covering number $\Omega(r^{-p})$, by a local nonsingular
coordinate chart. For the logarithmic consequence, integrate over a
small ball: on radii between $\sqrt\epsilon$ and a fixed radius the
radial integrand is bounded below by a constant times $r^{-1}$.
If there are at most $P_0$ phases per node and no other reductions,
$K\le(1+2nP_0)T$. Without this hypothesis the conclusion remains an
event lower bound.

\subsection{Aggregate loss at faces}
\label{app:propagation-face-loss}

\begin{proposition}[Two consequences of aggregate loss]
\label{prop:nd1}
Assume the objective-only event model above and the relaxation point test
with $\alpha>0$. Suppose on a set $N$,
$\ell(y)\ge\ell_0>0$, and on centered cubes of radius $\rho_1$ the
remainder in \eqref{prop:taylor-loss} is at most $K_1\delta^2$ when
all half-widths equal $\delta$. Put
$\rho_*\le\min\{\rho_1,\ell_0/(2K_1)\}$ and
$\vartheta=\ell_0\rho_*/2$.
Here $\rho_1>0$, $\rho_*>0$, and $\ell_0/(2K_1)=+\infty$ if $K_1=0$.

If $S\subseteq N\cap E(\eta)$ is a compact $C^1$ embedded arc whose unit
tangent satisfies $|\tau_i|\ge\tau_0>0$ for every coordinate, and
$\epsilon+\eta<\vartheta$, then
\begin{equation}
 K\ge\frac{\tau_0\mathcal H^1(S)}
 {2^{n+1}\sqrt{(\epsilon+\eta)/\alpha}
                 +8n(\epsilon+\eta)/\ell_0}.
 \label{prop:nd1-curve}
\end{equation}

Alternatively let $n\ge2$, $\rho_0>0$, $\gamma>0$, and
$N=Q_\infty(x_*,\rho_0)$, with
$Q_\infty(x_*,2\rho_0)\subseteq D$. Assume on the larger cube
$m\ge\gamma\|y-x_*\|_\infty^2$ and
$|\partial_im|\le\ell_0/8$. Take also $\rho_*\le\rho_0$. There is
a finite constant $C$, independent of $\epsilon$, such that
\begin{equation}
 K\ge C^{-1}\int_{N\cap\{m+\epsilon<\vartheta\}}
             (m+\epsilon)^{-n/2}\,dy.
 \label{prop:nd1-integral}
\end{equation}
If $m\le\Gamma\|y-x_*\|_\infty^2$ locally as well, this is
$\Omega(\log(1/\epsilon))$.
\end{proposition}

\begin{proof}
For a propagation-certified point choose
$\delta=\min\{\min_i(y_i-l_i,u_i-y_i),\rho_*\}$. The cube version
of \eqref{prop:taylor-loss} gives
$\delta<2g/\ell_0<\rho_*$, where $g=m+\epsilon$.
Thus one face is within $2g/\ell_0$.
Along the arc each tangent coordinate has constant sign, so each
coordinate is strictly monotone. A coordinate slab of half-width $r$
meets an arc of length at most $2r/\tau_0$. Relaxation-certified points
are in $2^n$ vertex balls of radius
$r_R=\sqrt{(\epsilon+\eta)/\alpha}$; propagation-certified points are in
$2n$ face slabs of half-width $r_F=2(\epsilon+\eta)/\ell_0$.
Their total arc length per event is at most
$(2^{n+1}r_R+4nr_F)/\tau_0$. Sum to obtain \eqref{prop:nd1-curve}.

For the integral, examine a face layer $y=(l_1+s,z)$ with
$0\le s<2g(y)/\ell_0<\rho_0$, and put $y'=(l_1,z)$.
The segment from $y'$ to $y$ lies in the larger cube. If
$M(z)=m(y')+\epsilon$, the gradient hypothesis gives
$|M-g(y)|<g(y)/4$. Therefore
\[
 s<\frac{8M(z)}{3\ell_0},\qquad
 g(y)^{-n/2}\le(5/4)^{n/2}M(z)^{-n/2}.
\]
Integrating first in $s$ bounds this layer by
\[
 \frac8{3\ell_0}(5/4)^{n/2}
 \int_{\|z-z_*\|_\infty\le2\rho_0}M(z)^{1-n/2}\,dz.
\]
For $n=2$ the last integrand is one. For $n>2$, use
$M\ge\gamma\|z-z_*\|_\infty^2$. In either case the integral is at most
$\gamma^{1-n/2}(n-1)2^n\rho_0$: the sup-norm radial integral is
$\int_0^{2\rho_0}t^{2-n}(n-1)2^{n-1}t^{n-2}\,dt$.
There are $2n$ layers. Relaxation-certified points contribute at most
$(\pi^2/(\alpha n))^{n/2}$ per event by the preceding arcsine estimate.
Adding these finite constants and summing the owned regions proves
\eqref{prop:nd1-integral}. The local upper quadratic bound makes its
integral logarithmic by the same radial estimate used above.
\end{proof}

\subsection{Two schedule-dependent round estimates}
\label{app:propagation-rounds}

\begin{proposition}[A slow expanded quadratic]
\label{prop:slow-rounds}
Let $a>0$, $0<\epsilon\le a^2/16$, and use the distinct term nodes
$s^2$ and $s$ in $f=s^2-2as+a^2$. A round consists of a forward pass,
cutoff $-\epsilon$, and a backward sum pass followed by the power revise.
Start from $[a-y_0,a+x_0]$, with $0<x_0,y_0\le a/4$ fixed. If
$z_0=\min\{x_0,y_0\}$, the box remains nonempty for at least
\[
 \frac a{4\sqrt\epsilon}
 \left(\arctan\frac{z_0}{\sqrt\epsilon}
             -\arctan\frac{4\sqrt\epsilon}{a}\right)-1
\]
rounds. As $\epsilon\downarrow0$ this is
$(\pi/8-o(1))a\epsilon^{-1/2}$. Its fixed point is nevertheless empty.
\end{proposition}

\begin{proof}
While the current endpoints straddle $a$, write them as $a-y,a+x$.
Backward propagation from the sum uses the old forward term ranges and
gives
\[
 y'=y-\frac{y^2+\epsilon}{2a},\qquad
 x'=\sqrt{a^2+2ax-\epsilon}-a
                  \ge x-\frac{x^2+\epsilon}{a}.
\]
The last inequality follows by squaring; its difference reduces to
$(x^2+\epsilon)(1+2x/a-(x^2+\epsilon)/a^2)\ge0$.
For $4\epsilon/a\le z\le a/4$, the decrement
$\Delta=(z^2+\epsilon)/a$ is at most $z/2$.
Both updated distances are therefore positive whenever both old distances
are at least $4\epsilon/a$, so that the round is nonempty. In this range
the potential $\theta(z)=\arctan(z/\sqrt\epsilon)$ falls by at most
\[
 \Delta\frac{\sqrt\epsilon}{\epsilon+(z-\Delta)^2}
 \le\frac{4\sqrt\epsilon}{a}.
\]
Apply this separately to $x$ and $y$ until their minimum first falls below
$4\epsilon/a$. The displayed number of rounds follows, with one round
allowed for crossing the threshold. The graph is one-sided on this
positive interval, so its limiting bound is zero and its fixed point
at cutoff $-\epsilon$ is empty.
\end{proof}

\begin{proposition}[Fast contraction at an exposed endpoint]
\label{prop:fast-rounds}
For $u=t^2$, $v=u^2$, $f=u-2v$, use forward-cutoff-backward rounds with
cutoff $-\epsilon$. If the initial upper bound on $u$ satisfies $0<U_0<1/2$,
the box empties in $O(\log\log(1/\epsilon))$ rounds as
$\epsilon\downarrow0$.
\end{proposition}

\begin{proof}
Each forward pass gives $v\le U_k^2$. The backward sum pass gives
$u\le2U_k^2-\epsilon$, hence $2U_{k+1}\le(2U_k)^2$ while nonempty.
Induction yields $2U_k\le(2U_0)^{2^k}$. The next pass is empty once
$(2U_0)^{2^{k+1}}<2\epsilon$, because $u\ge0$.
For $0<\epsilon<1/2$, taking
$k=1+\max\{0,\lceil\log_2(\log(1/(2\epsilon)) /
\log(1/(2U_0)))\rceil\}$ is more than sufficient and has the stated order.
\end{proof}
